\section{Fundamentação e trabalhos relacionados}
\label{sc:fundamentacao}

\subsection{Geração aumentada por recuperação}

Na formulação de Lewis et al.~\cite{lewis2020rag}, um sistema RAG combina um componente de recuperação, que seleciona trechos de uma coleção de documentos a partir da pergunta, e um componente gerador, que condiciona a resposta a esses trechos. Em implementações atuais, os documentos são divididos em fragmentos (\emph{chunks}) \cite{schwaber2025chunking}, convertidos em vetores por um modelo de embedding e armazenados em um banco vetorial \cite{ma2024vectorsurvey}; a pergunta é codificada pelo mesmo modelo, e os $k$ fragmentos mais próximos são entregues ao LLM. A arquitetura reduz alucinações e permite atualizar o conhecimento sem retreinar o modelo \cite{andriopoulos2023augmenting}, mas sua robustez depende de decisões tomadas empiricamente: Barnett et al.~\cite{barnett2024seven} catalogam sete pontos de falha recorrentes, entre eles a recuperação de trechos errados e a incapacidade do gerador de extrair a resposta de um contexto que a contém.

\subsection{Métodos de injeção de contexto}

Uma vez recuperados os fragmentos, há diferentes formas de entregá-los ao LLM \cite{topsakal2023langchain}. No \textbf{Stuff}, todos os fragmentos são concatenados em um único prompt. No \textbf{Refine}, a resposta é construída a partir do primeiro fragmento e refinada sequencialmente a cada fragmento seguinte. No \textbf{Map-Reduce}, cada fragmento gera uma resposta parcial independente, e uma chamada final combina as parciais. No \textbf{Map-Rerank}, cada fragmento gera uma resposta acompanhada de uma nota de confiança, e mantém-se apenas a de maior confiança. Şakar e Emekci~\cite{sakar2025rag} acrescentam dois métodos que atuam sobre a pergunta: o \textbf{Query Step-Down} gera perguntas alternativas mais específicas, responde cada uma e combina as respostas; o \textbf{Reciprocal RAG} segue o mesmo primeiro passo, mas retém apenas a resposta mais confiável. Ambos dialogam com a reescrita de consultas \cite{ma2023query} e com o RAG-Fusion \cite{rackauckas2024ragfusion}, que ampliam a cobertura da recuperação ao custo de chamadas adicionais ao LLM e do risco de desvio em relação à intenção original da pergunta.

A escolha do método tem efeito direto sobre o número de chamadas ao LLM e, portanto, sobre tokens e latência. Há também um efeito sobre a qualidade: Liu et al.~\cite{liu2023lost} mostram que LLMs aproveitam pior a informação posicionada no meio de contextos longos, o que penaliza estratégias que apenas acumulam fragmentos, e Nair et al.~\cite{nair2023drilling} mostram que representações condensadas de documentos longos preservam quase todo o desempenho com uma fração dos tokens.

\subsection{Compressão contextual e bancos vetoriais}

A compressão contextual reduz o volume de texto enviado ao LLM descartando fragmentos pouco relevantes. Em sua forma mais simples, adotada por Şakar e Emekci~\cite{sakar2025rag}, descartam-se os fragmentos cuja similaridade com a pergunta fica abaixo de um limiar. Os autores observaram economia de recursos, mas também queda de qualidade quando o limiar é punitivo demais para o domínio. Quanto à camada de armazenamento, bancos como o ChromaDB oferecem persistência em disco e filtros por metadados, enquanto bibliotecas como o FAISS \cite{johnson2021faiss} priorizam a velocidade da busca por similaridade em memória. Como ambos podem operar sobre os mesmos vetores com a mesma métrica de distância, a diferença entre eles é essencialmente de infraestrutura, e não de qualidade de recuperação.

\subsection{Avaliação de sistemas RAG}

A similaridade de cosseno entre os embeddings da resposta gerada e da resposta de referência é uma métrica barata e reprodutível, mas mede proximidade semântica global: penaliza respostas corretas redigidas de forma mais longa que a referência e não verifica se a resposta está ancorada nos documentos recuperados. O \emph{framework} RAGAS \cite{es2024ragas} propõe métricas específicas para RAG calculadas com um LLM como juiz: \emph{Faithfulness} (fração das afirmações da resposta sustentadas pelo contexto recuperado), \emph{Answer Relevancy} (quanto a resposta endereça a pergunta), \emph{Context Precision} e \emph{Context Recall} (qualidade e suficiência do contexto recuperado em relação à referência). O uso de LLMs como juízes, porém, tem vieses documentados, entre eles o de autopreferência: juízes tendem a avaliar melhor textos gerados por eles mesmos \cite{zheng2023judging, panickssery2024selfpref}. Isso é relevante quando, por restrição de orçamento, o mesmo modelo gera e avalia as respostas.

\subsection{RAG em português}

A avaliação de RAG em português ainda é incipiente. Medeiros e Oliveira~\cite{medeiros2025embeddings} compararam modelos de embedding e LLMs abertos e proprietários em três bases em português (Pirá, SQuAD 2 e FairytaleQA PT-BR), com ROUGE-L, BERTScore e as métricas RAGAS de fidelidade e relevância da resposta --- estas calculadas com o GPT-4o-mini como juiz e, por restrição orçamentária, apenas em duas das bases. Os melhores resultados foram obtidos pelo Multilingual E5 large e pelo Gemma 2 9B. Aquele estudo varia o embedding e o LLM com um método de injeção fixo; o presente trabalho faz o recorte complementar, fixando embedding e LLM e variando o método de injeção, a compressão contextual e o banco vetorial, sobre documentos longos de domínios profissionais.

\subsection{Lacuna}

A literatura avançou na acurácia \cite{ma2023query, rackauckas2024ragfusion}, na confiabilidade \cite{andriopoulos2023augmenting, barnett2024seven} e no entendimento das limitações de contexto \cite{liu2023lost, nair2023drilling} de sistemas RAG, e aplicações em domínios verticais --- jurídico \cite{cui2023chatlaw} e financeiro \cite{gupta2023gptinvestar} --- mostram a relevância prática do tema. A comparação sistemática entre métodos de injeção de contexto, porém, foi feita apenas em inglês \cite{sakar2025rag} e com uma única métrica de qualidade. Não há, até onde foi possível verificar, uma replicação desse desenho com documentos em português nem uma verificação de quanto as conclusões dependem da métrica escolhida. Este trabalho se insere nessa lacuna.
